% GENERATED FILE. Edit canonical lessons, then run npm run generate:books.
% canonical-source-hash: fnv1a64:238f8cf26abf40f4
% canonical-lessons: SA-C524-spardha, SA-C524-praksipati, SA-C524-vijeta, SA-C524-parajitah, SA-C524-praharati

\chapter{The Big Match}
\label{ch:sa-the-big-match}
\hlchaptermodality{524}

\begin{chapteropening}
\textbf{By the end of this chapter:} \emph{I can talk about a match and who won it.}

Complete the last lesson of chapter 524: I can talk about a match and who won it.
\end{chapteropening}

\section[\texorpdfstring{\sk{स्पर्धा} (spardhā)}{स्पर्धा (spardhā)}]{\sk{स्पर्धा} (spardhā) — a contest; rivalry}

\label{lesson:SA-C524-spardha}



\begin{quote}
Before the new one: say the Sanskrit for backstage, then the Sanskrit for acting.
\end{quote}

\subsection*{You'll want to know: \sk{स्पर्धा}}
\textbf{\sk{स्पर्धा}} — \emph{spardhā} — \textquotedblleft{}a contest; rivalry\textquotedblright{}.

\begin{quote}
\sk{श्वः} \sk{विद्यालये} \sk{क्रीडायाः} \sk{स्पर्धा} \sk{अस्ति।}
\end{quote}

\begin{quote}
\emph{śvaḥ vidyālaye krīḍāyāḥ spardhā asti.}
\end{quote}

\begin{quote}
Tomorrow there is a sports contest at school.
\end{quote}

A feminine noun.

\begin{cousinweb}
\textbf{\sk{स्पर्धा}} (\emph{spardhā}) comes from the root √spardh, \textquotedblleft{}to contend, to vie with\textquotedblright{}; MW gives \emph{spardhā} as emulation, rivalry, competition (MW).
\end{cousinweb}

\begin{grammarlens}[title={What you've built}]
One more word for describing what you lived and felt.
\end{grammarlens}

\subsection*{Your turn}
\begin{itemize}
\raggedright
  \item \emph{Say it:} \emph{spardhā}
  \item \emph{Say it:} \emph{spardhā}, once more
  \item \emph{Say it:} \emph{śvaḥ vidyālaye krīḍāyāḥ spardhā asti.} — Tomorrow there is a sports contest at school.
  \item \emph{Recall:} say the Sanskrit for backstage, then the Sanskrit for acting, then say \emph{spardhā} again
\end{itemize}

\begin{tcolorbox}[breakable,colback=teal!4,colframe=teal!35!black,title={Before you move on}]
What does \sk{स्पर्धा} mean? (\textquotedblleft{}A contest; rivalry\textquotedblright{}.) Say it once more.
\end{tcolorbox}
\section[\texorpdfstring{\sk{प्रक्षिपति} (prakṣipati)}{प्रक्षिपति (prakṣipati)}]{\sk{प्रक्षिपति} (prakṣipati) — hurls}

\label{lesson:SA-C524-praksipati}



\begin{quote}
Before the new one: say the Sanskrit for acting, then the Sanskrit for a contest.
\end{quote}

\subsection*{You'll want to know: \sk{प्रक्षिपति}}
\textbf{\sk{प्रक्षिपति}} — \emph{prakṣipati} — \textquotedblleft{}hurls\textquotedblright{}.

\begin{quote}
\sk{क्रीडकः} \sk{कन्दुकं} \sk{दूरं} \sk{प्रक्षिपति।}
\end{quote}

\begin{quote}
\emph{krīḍakaḥ kandukaṃ dūraṃ prakṣipati.}
\end{quote}

\begin{quote}
The player hurls the ball a long way.
\end{quote}

Past (he …): \sk{सः} \sk{कन्दुकं} \sk{प्रक्षिप्तवान्}, \emph{saḥ kandukaṃ prakṣiptavān}.

\begin{cousinweb}
\textbf{\sk{प्रक्षिपति}} (\emph{prakṣipati}) is built on \emph{pra}, \textquotedblleft{}forth\textquotedblright{}, and the root √kṣip, \textquotedblleft{}to throw\textquotedblright{}, which also gives \emph{kṣipati}; MW gives \emph{prakṣip} as to cast, to hurl, to fling (MW).
\end{cousinweb}

\begin{grammarlens}[title={What you've built}]
One more word for describing what you lived and felt.
\end{grammarlens}

\subsection*{Your turn}
\begin{itemize}
\raggedright
  \item \emph{Say it:} \emph{prakṣipati}
  \item \emph{Say it:} \emph{prakṣipati}, once more
  \item \emph{Say it:} \emph{krīḍakaḥ kandukaṃ dūraṃ prakṣipati.} — The player hurls the ball a long way.
  \item \emph{Recall:} say the Sanskrit for acting, then the Sanskrit for a contest, then say \emph{prakṣipati} again
\end{itemize}

\begin{tcolorbox}[breakable,colback=teal!4,colframe=teal!35!black,title={Before you move on}]
What does \sk{प्रक्षिपति} mean? (\textquotedblleft{}Hurls\textquotedblright{}.) Say it once more.
\end{tcolorbox}
\section[\texorpdfstring{\sk{विजेता} (vijetā)}{विजेता (vijetā)}]{\sk{विजेता} (vijetā) — a victor}

\label{lesson:SA-C524-vijeta}



\begin{quote}
Before the new one: say the Sanskrit for a contest, then the Sanskrit for hurls.
\end{quote}

\subsection*{You'll want to know: \sk{विजेता}}
\textbf{\sk{विजेता}} — \emph{vijetā} — \textquotedblleft{}a victor\textquotedblright{}.

\begin{quote}
\sk{एतस्यां} \sk{स्पर्धायाम्} \sk{अरुणः} \sk{विजेता} \sk{अस्ति।}
\end{quote}

\begin{quote}
\emph{etasyāṃ spardhāyām aruṇaḥ vijetā asti.}
\end{quote}

\begin{quote}
Arun is the winner of this contest.
\end{quote}

A masculine noun.

\begin{cousinweb}
\textbf{\sk{विजेता}} (\emph{vijetā}) is built on \emph{vi} and the root √ji, \textquotedblleft{}to conquer\textquotedblright{}, which also gives \emph{jayati}, with the agent ending \emph{-tṛ}; MW gives \emph{vijetṛ} as a vanquisher, a conqueror (MW).
\end{cousinweb}

\begin{grammarlens}[title={What you've built}]
One more word for describing what you lived and felt.
\end{grammarlens}

\subsection*{Your turn}
\begin{itemize}
\raggedright
  \item \emph{Say it:} \emph{vijetā}
  \item \emph{Say it:} \emph{vijetā}, once more
  \item \emph{Say it:} \emph{etasyāṃ spardhāyām aruṇaḥ vijetā asti.} — Arun is the winner of this contest.
  \item \emph{Recall:} say the Sanskrit for a contest, then the Sanskrit for hurls, then say \emph{vijetā} again
\end{itemize}

\begin{tcolorbox}[breakable,colback=teal!4,colframe=teal!35!black,title={Before you move on}]
What does \sk{विजेता} mean? (\textquotedblleft{}A victor\textquotedblright{}.) Say it once more.
\end{tcolorbox}
\section[\texorpdfstring{\sk{पराजितः} (parājitaḥ)}{पराजितः (parājitaḥ)}]{\sk{पराजितः} (parājitaḥ) — beaten; defeated}

\label{lesson:SA-C524-parajitah}



\begin{quote}
Before the new one: say the Sanskrit for hurls, then the Sanskrit for a victor.
\end{quote}

\subsection*{You'll want to know: \sk{पराजितः}}
\textbf{\sk{पराजितः}} — \emph{parājitaḥ} — \textquotedblleft{}beaten; defeated\textquotedblright{}.

\begin{quote}
\sk{मीरा} \sk{पराजिता} \sk{न} \sk{अस्ति।}
\end{quote}

\begin{quote}
\emph{mīrā parājitā na asti.}
\end{quote}

\begin{quote}
Mira is not beaten.
\end{quote}

\begin{cousinweb}
\textbf{\sk{पराजितः}} (\emph{parājitaḥ}) is built on \emph{parā}, \textquotedblleft{}away\textquotedblright{}, and the root √ji, \textquotedblleft{}to conquer\textquotedblright{}, as its past participle; MW gives \emph{parājita} as conquered, defeated, overthrown (MW).
\end{cousinweb}

\begin{grammarlens}[title={What you've built}]
One more word for describing what you lived and felt.
\end{grammarlens}

\subsection*{Your turn}
\begin{itemize}
\raggedright
  \item \emph{Say it:} \emph{parājitaḥ}
  \item \emph{Say it:} \emph{parājitaḥ}, once more
  \item \emph{Say it:} \emph{mīrā parājitā na asti.} — Mira is not beaten.
  \item \emph{Recall:} say the Sanskrit for hurls, then the Sanskrit for a victor, then say \emph{parājitaḥ} again
\end{itemize}

\begin{tcolorbox}[breakable,colback=teal!4,colframe=teal!35!black,title={Before you move on}]
What does \sk{पराजितः} mean? (\textquotedblleft{}Beaten; defeated\textquotedblright{}.) Say it once more.
\end{tcolorbox}
\section[\texorpdfstring{\sk{प्रहरति} (praharati)}{प्रहरति (praharati)}]{\sk{प्रहरति} (praharati) — hits}

\label{lesson:SA-C524-praharati}



\begin{quote}
Before the new one: say the Sanskrit for a victor, then the Sanskrit for beaten.
\end{quote}

\subsection*{You'll want to know: \sk{प्रहरति}}
\textbf{\sk{प्रहरति}} — \emph{praharati} — \textquotedblleft{}hits\textquotedblright{}.

\begin{quote}
\sk{क्रीडकः} \sk{कन्दुकं} \sk{दण्डेन} \sk{प्रहरति।}
\end{quote}

\begin{quote}
\emph{krīḍakaḥ kandukaṃ daṇḍena praharati.}
\end{quote}

\begin{quote}
The player hits the ball with a stick.
\end{quote}

Past (he …): \sk{सः} \sk{कन्दुकं} \sk{प्रहृतवान्}, \emph{saḥ kandukaṃ prahṛtavān}.

\begin{cousinweb}
\textbf{\sk{प्रहरति}} (\emph{praharati}) is built on \emph{pra}, \textquotedblleft{}forth\textquotedblright{}, and the root √hṛ, \textquotedblleft{}to carry\textquotedblright{}, which also gives \emph{harati}; MW gives \emph{prahṛ} as to strike, to hit (MW).
\end{cousinweb}

\begin{grammarlens}[title={What you've built}]
One more word for describing what you lived and felt.
\end{grammarlens}

\subsection*{Your turn}
\begin{itemize}
\raggedright
  \item \emph{Say it:} \emph{praharati}
  \item \emph{Say it:} \emph{praharati}, once more
  \item \emph{Say it:} \emph{krīḍakaḥ kandukaṃ daṇḍena praharati.} — The player hits the ball with a stick.
  \item \emph{Recall:} say the Sanskrit for a victor, then the Sanskrit for beaten, then say \emph{praharati} again
\end{itemize}

\begin{tcolorbox}[breakable,colback=teal!4,colframe=teal!35!black,title={Before you move on}]
What does \sk{प्रहरति} mean? (\textquotedblleft{}Hits\textquotedblright{}.) Say it once more.
\end{tcolorbox}
